\chapter{Positron Emission Tomography (PET) Scan}

\section*{What Is This Test?}
PET scans use a radioactive tracer (usually FDG, a sugar analog) to show metabolic activity. Cancer cells, infections, and inflammation consume more sugar and light up on PET. Often combined with CT (PET/CT) for precise localization.

\section*{Alternative Names}
PET scan, FDG-PET, PET/CT, positron emission tomography, whole-body PET


\section*{Principle}
PET uses a cyclotron-produced positron-emitting radiotracer, most commonly FDG (fluorodeoxyglucose) labeled with F-18. FDG accumulates in cells with high glucose metabolism (cancer cells, inflammation). Positrons annihilate with electrons, producing two 511 keV photons at 180 degrees, detected by coincidence circuitry in the PET scanner.

\section*{History}
PET scanning was developed in the 1970s by Michael Phelps and colleagues. The combination of PET with CT (PET/CT) in the early 2000s revolutionized oncologic imaging by merging metabolic and anatomic information.

\section*{Why This Test Is Ordered}
Your doctor orders a PET scan when they need to see how your cells are functioning at a molecular level. PET is most commonly used to detect cancer, determine whether cancer has spread, assess whether cancer treatment is working, and evaluate brain and heart function. It answers the question: are cells in this area of your body abnormally active, and could that be cancer, infection, or inflammation?

\section*{Is This a Primary or Secondary Test}
A secondary test. Usually performed after other imaging has identified a problem or as part of cancer staging.

\section*{How This Test Is Performed}
You receive an injection of a radioactive sugar (FDG) into your vein and wait about sixty minutes while the sugar distributes throughout your body. Cancer cells and other abnormally active cells absorb more sugar than normal cells, so they light up on the scan. You then lie on a table that passes through a PET/CT scanner, which combines functional information (from PET) with anatomical information (from CT). The scan takes thirty to sixty minutes. You should avoid strenuous exercise for 24 hours before the test and fast for four to six hours beforehand.

\section*{What Preparation Is Needed}
\begin{itemize}
\item Fast for four to six hours before the scan (water is allowed)
\item Avoid strenuous exercise for 24 hours before the test
\item If you are diabetic, your doctor will give you specific instructions about blood sugar management
\item Wear comfortable clothing without metal zippers or snaps
\item Tell your doctor if you are pregnant or breastfeeding
\end{itemize}

\section*{Contraindications and Precautions}
PET scans use a radioactive tracer and are generally not performed during pregnancy. Inform the team if you are pregnant or breastfeeding. The tracer passes into breast milk; discuss timing with your doctor. Blood sugar control is important for diabetic patients.

\section*{Risks}
PET scanning involves a small dose of radioactive tracer, with radiation exposure comparable to a CT scan. The tracer may cause a brief cold sensation at the injection site. Rare allergic reactions may occur. Inform the team if you are pregnant or breastfeeding.

\section*{What Happens After the Test}
You can resume normal activities immediately. You should drink plenty of fluids to help flush the radioactive tracer from your body. Results typically take one to twenty-four hours as the images require detailed interpretation by a radiologist.

\section*{Understanding Your Results}

\section*{Normal Results}
A normal result shows no areas of abnormally increased metabolic activity that would suggest cancer, inflammation, or infection.

\section*{Follow-up Tests}
If the PET scan shows abnormal metabolic activity, your doctor may order a biopsy for tissue diagnosis, a CT scan for more detailed anatomical evaluation, an MRI for soft tissue characterization, or follow-up PET scanning to assess treatment response.
\clearpage
